\documentclass[10pt,a4paper]{amsart}
\usepackage{amsmath,amssymb,amsthm,graphicx}
% For commutative diagrams we use Xy-pic
\usepackage[all]{xy}

% We use 2cell for 2-commutative diagrams.
\xyoption{2cell}
% Theorem environments.
%
\theoremstyle{plain}
\newtheorem{theorem}[subsection]{Theorem}
\newtheorem{proposition}[subsection]{Proposition}
\newtheorem{lemma}[subsection]{Lemma}

\theoremstyle{definition}
\newtheorem{definition}[subsection]{Definition}
\newtheorem{example}[subsection]{Example}
\newtheorem{exercise}[subsection]{Exercise}
\newtheorem{situation}[subsection]{Situation}

\theoremstyle{remark}
\newtheorem{remark}[subsection]{Remark}
\newtheorem{remarks}[subsection]{Remarks}

\numberwithin{equation}{subsection}

% Macros
%
\def\lim{\mathop{\mathrm{lim}}\nolimits}
\def\colim{\mathop{\mathrm{colim}}\nolimits}
\def\Spec{\mathop{\mathrm{Spec}}}
\def\Hom{\mathop{\mathrm{Hom}}\nolimits}
\def\Ext{\mathop{\mathrm{Ext}}\nolimits}
\def\SheafHom{\mathop{\mathcal{H}\!\mathit{om}}\nolimits}
\def\SheafExt{\mathop{\mathcal{E}\!\mathit{xt}}\nolimits}
\def\Sch{\mathit{Sch}}
\def\Mor{\mathop{\mathrm{Mor}}\nolimits}
\def\Ob{\mathop{\mathrm{Ob}}\nolimits}
\def\Sh{\mathop{\mathit{Sh}}\nolimits}
\def\NL{\mathop{N\!L}\nolimits}
\def\CH{\mathop{\mathrm{CH}}\nolimits}
\def\proetale{{pro\text{-}\acute{e}tale}}
\def\etale{{\acute{e}tale}}
\def\QCoh{\mathit{QCoh}}
\def\Ker{\mathop{\mathrm{Ker}}}
\def\Im{\mathop{\mathrm{Im}}}
\def\Coker{\mathop{\mathrm{Coker}}}
\def\Coim{\mathop{\mathrm{Coim}}}

% Boxtimes
%
\DeclareMathSymbol{\boxtimes}{\mathbin}{AMSa}{"02}

%
% Macros for moduli stacks/spaces
%
\def\QCohstack{\mathcal{QC}\!\mathit{oh}}
\def\Cohstack{\mathcal{C}\!\mathit{oh}}
\def\Spacesstack{\mathcal{S}\!\mathit{paces}}
\def\Quotfunctor{\mathrm{Quot}}
\def\Hilbfunctor{\mathrm{Hilb}}
\def\Curvesstack{\mathcal{C}\!\mathit{urves}}
\def\Polarizedstack{\mathcal{P}\!\mathit{olarized}}
\def\Complexesstack{\mathcal{C}\!\mathit{omplexes}}
% \Pic is the operator that assigns to X its picard group, usage \Pic(X)
% \Picardstack_{X/B} denotes the Picard stack of X over B
% \Picardfunctor_{X/B} denotes the Picard functor of X over B
\def\Pic{\mathop{\mathrm{Pic}}\nolimits}
\def\Picardstack{\mathcal{P}\!\mathit{ic}}
\def\Picardfunctor{\mathrm{Pic}}
\def\Deformationcategory{\mathcal{D}\!\mathit{ef}}
\usepackage{cleveref}
\crefname{originalreference}{original reference}{original references}
\Crefname{originalreference}{Original reference}{Original references}
\makeatletter
\newlabel{lemma-colimit-constructible}{{095P}{source}{}{sourcepointer-1}{}}
\newlabel{lemma-colimit-constructible@cref}{{[originalreference][0][]095P}{[1][0][]source}{}{}{}}
\makeatother
\crefname{originallemma}{original lemma}{original lemmas}
\Crefname{originallemma}{Original lemma}{Original lemmas}
\makeatletter
\newlabel{sites-modules-lemma-module-quotient-direct-sum}{{093B}{source}{}{sourcepointer-2}{}}
\newlabel{sites-modules-lemma-module-quotient-direct-sum@cref}{{[originallemma][0][]093B}{[1][0][]source}{}{}{}}
\makeatother
\title{A constructible cohomology sheaf of a complex is hit by a finite sum of extensions by zero of the constant sheaf}
\author{The Stacks Project Authors}
\date{stacks-project:etale-cohomology.tex:local-snapshot-20261011; GNU FDL 1.2 or later}
\begin{document}
\maketitle
\noindent\textit{Editorial note.} All delivered mathematical text is the author's own, copied from the pinned Stacks Project chapter file etale-cohomology.tex; no statement or proof step was rewritten, reordered or removed.\par\medskip
\noindent\textit{Editorial note.} The counted proof is the author's complete argument: the proof is written out in full: the projection of the kernel onto the cohomology sheaf, the surjection from extension sheaves over affine objects, and the finiteness step from constructibility.\par\medskip
\noindent\textit{Editorial note.} The label lemma-one-constructible names the counted statement, and the proof environment immediately following it in the source is delivered verbatim as its proof.\par\medskip
\noindent\textit{Editorial note.} The counted statement is carried by the author's own lemma environment in the chapter source.\par\medskip
\noindent\textit{Editorial note.} The statement and proof are copied from the pinned chapter file; the two cited inputs are kept as external pointers with their canonical Stacks tags.\par\medskip
\noindent\textit{Allowed standard prerequisites (editorial).}\begin{itemize}
\item Surjection onto a module from a direct sum of extensions by zero. Modules on Sites, the lemma representing a sheaf of modules as a quotient of a direct sum of j\_\{U!\} of the structure sheaf over affine objects. Supplies the surjection onto the kernel of K\textasciicircum{}b -> K\textasciicircum{}\{b+1\} used at the start.
\item Constructible sheaves are finitely generated over the colimit. Same chapter, the lemma stating that a constructible sheaf is the union of the images of maps from finitely many extension sheaves. Supplies that the image of a finite sub-sum already equals H\textasciicircum{}b(K).
\end{itemize}\medskip
\ifdefined\hypertarget\hypertarget{sourcepointer-1}{}\fi
\noindent\textit{Source reference.} Original reference 095P: sources/stacks-etale-cohomology/pinned-original-source.tex, source lines 11903--11903. Prerequisite pointer invoked at proof line(s) 13637; it names a standard fact of the same chapter that the counted argument uses, and every step of the counted proof built on it is retained verbatim in the delivered body.\par\medskip
\ifdefined\hypertarget\hypertarget{sourcepointer-2}{}\fi
\noindent\textit{Source reference.} Original lemma 093B: sources/stacks-chow/tags.txt, source lines 11509--11509. Prerequisite pointer invoked at proof line(s) 13630; it names a standard fact of a sibling chapter that the counted argument uses, and every step of the counted proof built on it is retained verbatim in the delivered body.\par\medskip
\begin{lemma}
\label{lemma-one-constructible}
Let $X$ be a quasi-compact and quasi-separated scheme.
Let $\Lambda$ be a Noetherian ring. Let $K \in D(X_\etale, \Lambda)$
and $b \in \mathbf{Z}$ such that $H^b(K)$ is constructible.
Then there exist a sheaf $\mathcal{F}$ which is a finite direct sum
of $j_{U!}\underline{\Lambda}$ with $U \in \Ob(X_\etale)$ affine and
a map $\mathcal{F}[-b] \to K$ in $D(X_\etale, \Lambda)$
inducing a surjection $\mathcal{F} \to H^b(K)$.
\end{lemma}

\begin{proof}
Represent $K$ by a complex $\mathcal{K}^\bullet$ of sheaves of
$\Lambda$-modules. Consider the surjection
$$
\Ker(\mathcal{K}^b \to \mathcal{K}^{b + 1})
\longrightarrow
H^b(K)
$$
By Modules on Sites, Lemma
\ref{sites-modules-lemma-module-quotient-direct-sum}
we may choose a surjection
$\bigoplus_{i \in I} j_{U_i!} \underline{\Lambda} \to
\Ker(\mathcal{K}^b \to \mathcal{K}^{b + 1})$
with $U_i$ affine. For $I' \subset I$ finite, denote
$\mathcal{H}_{I'} \subset H^b(K)$ the image of
$\bigoplus_{i \in I'} j_{U_i!} \underline{\Lambda}$. By
Lemma \ref{lemma-colimit-constructible}
we see that $\mathcal{H}_{I'} = H^b(K)$ for some $I' \subset I$ finite.
The lemma follows taking
$\mathcal{F} = \bigoplus_{i \in I'} j_{U_i!} \underline{\Lambda}$.
\end{proof}

\end{document}
